\section{Conclusion}

Nibli makes several decisions visible that an ordinary Boolean query interface
can conceal: whether a result is definitive, which premises support it, which
closed-world assumptions it uses, and which evaluation profile produced it.
The evaluated design combines bounded reasoning, selective completion,
source-preserving withdrawal, and an evidence interface across shared native
and WebAssembly code.

The experiments record no definitive disagreement or stale source citation,
and the shared runtime scenarios agree. Those observations coexist with
substantial limits: single-statement ingestion fails to finish at the largest
policy target, forward recursion hits the cutoff on short chains, and
certificate-producing queries cost more than verdict queries. Selective
materialization improves returned definitiveness in the forward-chain slice.
These results support the specific semantic, transition, and portability
claims recorded in the artifact, together with their performance boundaries.
The proof argument remains conditional at its model boundaries, and envelope
coherence remains distinct from semantic verification. Future extensions that
change materialization, witness handling, or persistence should preserve those
distinctions and rerun the affected evidence, rather than inherit claims from
an earlier revision.
